\documentclass[12pt]{article}
\usepackage[utf8]{inputenc}
\usepackage[T1]{fontenc}
\usepackage{hyperref}
\title{Teke-Tsaayi Language Resources for K-12 Curriculum Development}
\author{Compiled by Claude Sonnet 5 (claude-sonnet-5) for Andr\'as Kornai}
\date{}
\begin{document}
\maketitle

\section{Introduction}
Teke-Tsaayi (ISO 639-3 tyi, also called West Teke) is a Bantu language of the Republic of Congo. It belongs to a dialect continuum whose varieties include Tsaayi, Laali, Yaa and Tyee, and Tsaayi is the dominant variety. Speaker estimates in the sources found range from about 96,000 (ASJP) to about 120,000 (Omniglot). Joshua Project reports larger figures for the Tsaayi people group. This report is based on a small number of web searches. Very little material exists, and none of it is school material.

\section{Primary Educational Resources}
No primers, textbooks or school programmes for Teke-Tsaayi turned up in the search results.

\section{Linguistic and Reference Materials}
\begin{itemize}
\item A Teke Tyee--French dictionary (SIL International, 2022, edited by Chrisnah Renaudot Mfouhou and Ruth Raharimanantsoa) on Webonary. It covers the Tyee variety, which is a neighbouring variety of West Teke and not Tsaayi proper. It uses a standardized orthography: \url{https://www.webonary.work/teke-tyee/overview/foreword/?lang=en}
\item Omniglot page on West Teke (Teke-Tsaayi), which describes the Latin-based writing and links to alphabet material: \url{https://www.omniglot.com/writing/westteke.htm}
\item Wikipedia overview of West Teke: \url{https://en.wikipedia.org/wiki/West_Teke_language}
\item ASJP word-list entry for Western Teke Tsaayi: \url{https://asjp.clld.org/languages/WESTERN_TEKE_TSAAYI}
\item Glottolog entry (teke1281): \url{https://glottolog.org/resource/languoid/id/teke1281.bigmap.html}
\end{itemize}

\section{Religious and Community Resources}
Joshua Project lists Bible translation as started, with no portions reported: \url{https://www.joshuaproject.net/index.php/languages/tyi}. Its page on the Tsaayi in Congo is \url{https://www.joshuaproject.net/index.php/people_groups/15385/CF}. I found no community organisations.

\section{Digital and Audio Resources}
The Global Recordings Network page for Teke-Tsaayi lists no recordings as available: \url{https://globalrecordings.net/en/language/tyi}

\section{Curriculum Development Context}
There is no evidence of Teke-Tsaayi being taught in schools. The only substantial lexical resource is the Tyee dictionary, whose standardized orthography could serve as a starting point. A curriculum would need to build primers, texts and a Tsaayi-specific lexicon from scratch, ideally with community and SIL involvement. No Tsaayi-specific dictionary or grammar was found.
\end{document}